\section*{Capítulo \ref{ch:b2:metamorphic-facies} --- Reacciones y facies metamórficas}

\begin{solution}{exo:b2:metamorphic-facies:1}
(a) \omterm{def:b2:metamorphic-facies:metamorphism}{Metamorfismo} de contacto, un \omterm{def:b2:metamorphic-facies:metamorphism}{metamorfismo} térmico local debido al calor de la intrusión.
(b) \omterm{def:b2:metamorphic-facies:metamorphism}{Metamorfismo} de dislocación, en una zona de falla: reducción del tamaño de grano, una milonita.
(c) \omterm{def:b2:metamorphic-facies:metamorphism}{Metamorfismo} de fondo oceánico, por el agua de mar caliente cerca de la dorsal, con una
recristalización incompleta. (d) \omterm{def:b2:metamorphic-facies:metamorphism}{Metamorfismo} de enterramiento, regional, sin deformación ni
esquistosidad.
\end{solution}

\begin{solution}{exo:b2:metamorphic-facies:2}
$\Delta_r S = 95.40 - 82.80 = \qty{12.60}{J.K^{-1}}$, $\Delta_r V = 4.986 - 4.415 =
\qty{0.571}{J.bar^{-1}}$; pendiente $12.60/0.571 = \qty{22.1}{bar/K} = \qty{2.21}{MPa/K}$.
En \qty{100}{\celsius} la presión de equilibrio sube \qty{2.21}{kbar}.
\end{solution}

\begin{solution}{exo:b2:metamorphic-facies:3}
(a) Esquistos verdes. (b) Esquistos azules. (c) Anfibolitas: hornblenda con plagioclasa de más
de un \qty{17}{\percent} de anortita. (d) Eclogitas.
\end{solution}

\begin{solution}{exo:b2:metamorphic-facies:4}
(a) $\rho = M/V$: cianita $162.046/44.15 = \qty{3.670}{g.cm^{-3}}$, andalucita
$162.046/51.52 = 3.145$, sillimanita $162.046/49.86 = 3.250$. (b) La cianita: a mayor
presión el término $V\,(P - 1)$ de la energía de Gibbs crece menos para el menor volumen.
(c) Los cristales naturales contienen otros elementos (el hierro férrico, por ejemplo, sustituye a parte del
aluminio en la andalucita y la sillimanita naturales) y las medidas sobre ejemplares
naturales son menos precisas que los volúmenes molares de las tablas.
\end{solution}

\begin{solution}{exo:b2:metamorphic-facies:5}
(a) $\Delta_r G = 3900 - 298.15\times8.59 = \qty{+1339}{J} > 0$: cianita. (b) A
\qty{1}{bar}, $\Delta_r H = T\Delta_r S$: $T = 3900/8.59 = \qty{454}{K} =
\qty{181}{\celsius}$. (c) $P = (8.59\times673.15 - 3900)/0.737 = 1882/0.737 =
\qty{2554}{bar} = \qty{2.55}{kbar}$, frente a 2.54 medidos. (d) La reacción es demasiado lenta
a baja temperatura para avanzar: la andalucita sobrevive de forma metaestable, como la mayoría de los minerales
de alta temperatura en la superficie.
\end{solution}

\begin{solution}{exo:b2:metamorphic-facies:6}
(a) $\Delta_r H = 7700 + 1000\times(121.39 - 122.29) = \qty{6800}{J}$; $\Delta_r S = 296.44
- 285.02 = \qty{11.42}{J.K^{-1}}$; $\Delta_r G = 6800 - 1000\times11.42 = \qty{-4620}{J}$.
(b) Con los valores a \qty{298.15}{K}: $7700 - 1000\times12.60 = \qty{-4900}{J}$.
(c) Sillimanita. El método simple se equivoca en \qty{280}{J}, cerca de un \qty{6}{\percent}; dividido
por $\Delta_r V = 0.571$, son unos \qty{490}{bar} sobre la línea, mucho menos que el
efecto de la incertidumbre de las entalpías.
\end{solution}

\begin{solution}{exo:b2:metamorphic-facies:7}
(a) $\Delta_r G = -1127.4 + 1128.5 = \qty{1.1}{kJ}$ y $\Delta_r V = (34.15 -
36.93)/10 = \qty{-0.278}{J.bar^{-1}}$: $P = 1100/0.278 = \qty{3957}{bar} \approx
\qty{4.0}{kbar}$. (b) $\Delta_r S = 88.0 - 91.7 = -3.7$; pendiente $-3.7/-0.278 =
\qty{13.3}{bar/K}$. (c) Con $\Delta_r H = 0$, $P = T\Delta_r S/\Delta_r V = 13.31\,T$:
a \qty{573.15}{K}, \qty{7629}{bar} $= \qty{7.63}{kbar}$. (d) El aragonito de las conchas es
metaestable en la superficie: los organismos lo fabrican, y persiste porque solo se transforma en
calcita lentamente.
\end{solution}

\begin{solution}{exo:b2:metamorphic-facies:8}
(a) $-32.4/-1.698 = \qty{19.1}{bar/K}$. (b) $P = (T\Delta_r S - \Delta_r H)/\Delta_r V =
(-32.4\times773.15 + 5000)/(-1.698) = (-25050 + 5000)/(-1.698) = \qty{11808}{bar}
\approx \qty{11.8}{kbar}$. (c) $\sqrt{2.6^2 + 3.6^2 + 1.0^2} = \qty{4.55}{kJ}$ sobre
$\Delta_r H$, luego $4550/1.698 = \qty{2.7}{kbar}$ sobre la presión. (d) A las presiones de
la facies de las eclogitas el componente albita de la plagioclasa es inestable frente a la
jadeíta y el cuarzo; el sodio pasa al piroxeno (onfacita), y la plagioclasa
desaparece.
\end{solution}

\begin{solution}{exo:b2:metamorphic-facies:9}
(a) De 1 a \qty{2}{kbar}: $1000/34 = \qty{29.4}{bar/K}$; de 2 a 3: $1000/25 =
\qty{40.0}{bar/K}$. (b) Con $\Delta_r S$ fijo, $\Delta_r V$ es inversamente proporcional a
la pendiente: disminuyó en un factor $40.0/29.4 = 1.36$. (c) El agua liberada: su volumen
molar disminuye al aumentar la presión, mientras que los volúmenes de los sólidos apenas cambian.
\end{solution}

\begin{solution}{exo:b2:metamorphic-facies:10}
(a) A \qty{873.15}{K}: cianita $=$ andalucita $(8.59\times873.15 - 3900)/0.737 =
\qty{4885}{bar}$; cianita $=$ sillimanita $(12.60\times873.15 - 7700)/0.571 =
\qty{5782}{bar}$; andalucita $=$ sillimanita $(4.01\times873.15 - 3800)/(-0.166) =
\qty{1799}{bar}$. (b) A \qty{4885}{bar}, por debajo de la línea cianita $=$ sillimanita (5782),
la sillimanita es más estable que la cianita, y por encima de la línea andalucita $=$ sillimanita
(1799), más estable que la andalucita: el punto está en el campo de la sillimanita, y el
equilibrio cianita $=$ andalucita es allí metaestable. (c) A \qty{4}{kbar}: entre 1799
y \qty{5782}{bar}, sillimanita. A \qty{7}{kbar}: por encima de \qty{5782}{bar} la cianita gana a la
sillimanita, y por encima de \qty{4885}{bar} gana a la andalucita: cianita.
\end{solution}

\begin{solution}{exo:b2:metamorphic-facies:11}
(a) $z = P/(\rho g) = 3.75\times10^8/(2700\times9.81) = \qty{14.2}{km}$, con una incertidumbre de
$0.25\times10^8/26487 = \qty{0.94}{km}$. (b) $(504 - 10)/14.16 =
\qty{35}{\celsius.km^{-1}}$. (c) Los tres minerales coexisten en equilibrio solo en el
\omterm{def:b2:metamorphic-facies:invariant}{punto invariante}: la trayectoria de la roca debe pasar por él, y las reacciones deben
detenerse allí. La mayoría de las rocas pasan a su lado y conservan uno o dos polimorfos.
\end{solution}

\begin{solution}{exo:b2:metamorphic-facies:12}
(a) X: $T = 10 + 10\times30 = \qty{310}{\celsius}$ (\qty{583.15}{K}); $P = 3000\times
9.81\times30000 = \qty{8.83e8}{Pa} = \qty{8.83}{kbar}$. Y: $T = 10 + 30\times20 =
\qty{610}{\celsius}$ (\qty{883.15}{K}); $P = 2800\times9.81\times20000 = \qty{5.49}{kbar}$.
(b) X: calcita $=$ aragonito a $13.31\times583.15 = \qty{7.76}{kbar} < 8.83$, aragonito;
albita $=$ jadeíta + cuarzo a $(32.4\times583.15 - 5000)/1.698 = \qty{8.18}{kbar} <
8.83$, jadeíta + cuarzo; cianita $=$ andalucita a $(8.59\times583.15 - 3900)/0.737 =
\qty{1.51}{kbar}$, por debajo de la temperatura del punto triple: cianita. Y: calcita $=$
aragonito a \qty{11.75}{kbar} $> 5.49$, calcita; albita $=$ jadeíta + cuarzo a
\qty{13.9}{kbar}, albita; por encima del punto triple, cianita $=$ sillimanita a
\qty{6.00}{kbar} $> 5.49$ y andalucita $=$ sillimanita a \qty{1.56}{kbar} $< 5.49$:
sillimanita. (c) X, a \qty{10}{\celsius\per\km}, es la trayectoria más fría: la jadeíta y el
aragonito a solo \qty{310}{\celsius} la sitúan en la facies de los esquistos azules. Y, tres veces
más cálida, termina con sillimanita a \qty{610}{\celsius} y presión moderada, en la
facies de las anfibolitas.
\end{solution}

\begin{solution}{pb:b2:metamorphic-facies:1}
\noindent\textbf{Parte I.}

\textbf{1.} $\Delta_r H = -2589.9 + 2593.8 = \qty{3.9}{kJ}$; $\Delta_r S = 91.39 - 82.80 =
\qty{8.59}{J.K^{-1}}$; $\Delta_r V = 5.152 - 4.415 = \qty{0.737}{J.bar^{-1}}$.

\textbf{2.} $\qty{7.7}{kJ}$; $\qty{12.60}{J.K^{-1}}$; $\qty{0.571}{J.bar^{-1}}$.

\textbf{3.} $\qty{3.8}{kJ}$; $\qty{4.01}{J.K^{-1}}$; $\qty{-0.166}{J.bar^{-1}}$; cada uno es
el segundo menos el primero, ya que la andalucita se cancela.

\textbf{4.} La sillimanita tiene la mayor entropía (95.40): está del lado de alta temperatura
de sus dos líneas. La cianita tiene el menor volumen (4.415): está del lado de
alta presión de sus dos líneas. La andalucita, de mayor volumen, ocupa el
lado de baja presión.

\textbf{5.} $8.59/0.737 = \qty{11.66}{bar/K}$; $12.60/0.571 = \qty{22.07}{bar/K}$;
$4.01/(-0.166) = \qty{-24.16}{bar/K}$.

\textbf{6.} La sillimanita tiene más entropía pero menos volumen que la andalucita: $\Delta_r S > 0$
y $\Delta_r V < 0$. Subir la temperatura favorece a la sillimanita, y subir la presión
también, de modo que la línea va hacia presiones menores a temperaturas mayores.

\noindent\textbf{Parte II.}

\textbf{7.} $3900 - 298.15\times8.59 = \qty{+1339}{J}$; $7700 - 298.15\times12.60 =
\qty{+3943}{J}$. Ambos positivos: la cianita es estable en la superficie.

\textbf{8.} $T = \Delta_r H/\Delta_r S$: $3900/8.59 = \qty{454.0}{K} = \qty{181}{\celsius}$;
$7700/12.60 = \qty{611.1}{K} = \qty{338}{\celsius}$; $3800/4.01 = \qty{947.6}{K} =
\qty{674}{\celsius}$.

\textbf{9.} Al calentar a \qty{1}{bar}, la cianita cede el paso a la andalucita a
\qty{181}{\celsius}; a \qty{338}{\celsius} la cianita ya no es estable, así que cianita $=$
sillimanita es allí metaestable; la andalucita cede el paso a la sillimanita a
\qty{674}{\celsius}. La primera y la tercera son equilibrios estables.

\textbf{10.} $(8.59\times673.15 - 3900)/0.737 = \qty{2554}{bar} = \qty{2.55}{kbar}$,
\qty{0.01}{kbar} de los experimentos.

\textbf{11.} $(12.60\times973.15 - 7700)/0.571 = 4562/0.571 = \qty{7989}{bar} =
\qty{7.99}{kbar}$, \qty{0.13}{kbar} por encima de los experimentos.

\textbf{12.} Cianita $=$ andalucita sube desde \qty{181}{\celsius} a \qty{1}{bar} con
pendiente \qty{11.7}{bar/K}; andalucita $=$ sillimanita baja desde unos \qty{11.5}{kbar} a
\qty{200}{\celsius} hasta \qty{1}{bar} a \qty{674}{\celsius}; cianita $=$ sillimanita, la
más empinada de las dos líneas ascendentes, corta a ambas cerca de \qty{514}{\celsius} y \qty{3.9}{kbar}, y cada línea es estable a
un solo lado del cruce (el diagrama del capítulo).

\noindent\textbf{Parte III.}

\textbf{13.} $a_1 = 8.59/0.737 = \qty{11.655}{bar/K}$, $b_1 = 3900/0.737 =
\qty{5291.7}{bar}$; $a_2 = 12.60/0.571 = \qty{22.067}{bar/K}$, $b_2 = 7700/0.571 =
\qty{13485.1}{bar}$.

\textbf{14.} $T = (b_2 - b_1)/(a_2 - a_1) = 8193.4/10.412 = \qty{786.92}{K} =
\qty{513.8}{\celsius}$.

\textbf{15.} $P = 11.655\times786.92 - 5291.7 = \qty{3880}{bar} = \qty{3.88}{kbar}$.

\textbf{16.} $(4.01\times786.92 - 3800)/(-0.166) = (3155.5 - 3800)/(-0.166) =
\qty{3882}{bar}$: la misma presión salvo el redondeo de los datos, como exige la proposición del capítulo.

\textbf{17.} \qty{514}{\celsius} queda a \qty{10}{\celsius} de \qty{504}{\celsius}, dentro de
la incertidumbre de \qty{20}{\celsius}; \qty{3.88}{kbar} queda a \qty{0.13}{kbar} de 3.75,
dentro de 0.25. El acuerdo es bueno.

\textbf{18.} No: \qty{536}{\celsius} queda a \qty{32}{\celsius} de los experimentos, fuera de su
incertidumbre. Despreciar las capacidades caloríficas no era el error principal: lo son las entalpías.

\noindent\textbf{Parte IV.}

\textbf{19.} El $\Delta_r H$ de cianita $\to$ andalucita sube \qty{1000}{J}; a
$T$ fija, $P = (T\Delta_r S - \Delta_r H)/\Delta_r V$ baja $1000/0.737 = \qty{1357}{bar}$.

\textbf{20.} $\Delta T = -1357/(a_2 - a_1) = -1357/10.41 = \qty{-130}{K}$: el punto triple
se desplaza a unos \qty{384}{\celsius}.

\textbf{21.} La línea cianita $=$ sillimanita baja $1000/0.571 = \qty{1751}{bar}$, y
el punto triple se desplaza $+1751/10.41 = \qty{+168}{K}$.

\textbf{22.} Cada entalpía por separado, variada en $\pm\qty{2.0}{kJ}$, desplaza el punto triple
$\pm\qty{260}{K}$ (andalucita) o $\pm\qty{336}{K}$ (sillimanita): la calorimetría sola admite
puntos triples repartidos en varios cientos de grados.

\textbf{23.} Las entalpías tabuladas no son puramente calorimétricas: dentro de sus
incertidumbres se ajustan para que las líneas calculadas se ajusten a los experimentos de inversión.
Los experimentos fijan el diagrama; las tablas lo trasladan, de forma coherente, a otras reacciones.

\textbf{24.} El punto triple calculado a partir de las tablas con propiedades constantes está a
\textbf{\qty[text-series-to-math=true]{514}{\celsius}}.
\end{solution}
